\documentclass{eleclab-report}

\usetikzlibrary{arrows.meta}
\usepackage[siunitx, RPvoltages]{circuitikz}

\setexperimentnumber{6}
\setexperimenttitle{移相器的设计与测试}

\newcommand{\ExpSixTableSetup}{%
  \setlength{\tabcolsep}{2.8pt}%
  \renewcommand{\arraystretch}{1.18}%
}
\newcommand{\ExpSixTableFont}{\ElecLabBodyFont\fontsize{9pt}{12pt}\selectfont}
\newcommand{\ExpSixSmallTableFont}{\ElecLabBodyFont\fontsize{8.5pt}{11.5pt}\selectfont}
\newcommand{\ExpSixTable}[1]{%
  \begingroup
  \ExpSixTableSetup
  \begin{center}
  {\ExpSixTableFont #1}
  \end{center}
  \endgroup
}
\newcommand{\ExpSixSmallTable}[1]{%
  \begingroup
  \setlength{\tabcolsep}{2.0pt}%
  \renewcommand{\arraystretch}{1.18}%
  \begin{center}
  {\ExpSixSmallTableFont #1}
  \end{center}
  \endgroup
}
\newcommand{\AnalysisParagraph}[1]{\par\hspace*{2em}#1\par}
\newcommand{\MeasureBlank}{\LabTableBlank}

\newcommand{\PhaseShiftCircuit}{%
  \begin{circuitikz}[american,x=1cm,y=1cm,line cap=round,line join=round]
    \coordinate (a) at (0,2.6);
    \coordinate (b) at (0,0);
    \coordinate (c) at (5.6,2.6);
    \coordinate (d) at (5.6,0);
    \draw (a) node[ocirc,label=left:$a$] {} to[R,l=$R$] (c) node[ocirc,label=right:$c$] {};
    \draw (b) node[ocirc,label=left:$b$] {} to[R,l_=$R$] (d) node[ocirc,label=right:$d$] {};
    \draw (a) to[C,l_=$C$] (d);
    \draw (b) to[C,l=$C$] (c);
    \draw[-{Latex[length=2.0mm]}] (-0.45,2.25) -- (-0.45,0.35);
    \draw[-{Latex[length=2.0mm]}] (6.05,2.25) -- (6.05,0.35);
    \node[left] at (-0.5,1.3) {$U_1$};
    \node[right] at (6.1,1.3) {$U_2$};
    \node[anchor=base] at (2.8,-0.68) {\labminorfont X型RC移相电路};
  \end{circuitikz}%
}

\newcommand{\EquivalentPhaseCircuit}{%
  \begin{circuitikz}[american,x=1cm,y=1cm,line cap=round,line join=round]
    \coordinate (a) at (0,2.9);
    \coordinate (b) at (0,0);
    \coordinate (c) at (2.6,1.45);
    \coordinate (d) at (5.2,1.45);
    \draw (a) node[ocirc,label=left:$a$] {} to[R,l=$R$] (c) to[C,l=$C$] (b) node[ocirc,label=left:$b$] {};
    \draw (a) -- (5.2,2.9) to[C,l=$C$] (d) to[R,l=$R$] (5.2,0) -- (b);
    \draw (c) node[ocirc,label=below left:$c$] {};
    \draw (d) node[ocirc,label=below right:$d$] {};
    \draw[-{Latex[length=2.0mm]}] (-0.45,2.55) -- (-0.45,0.35);
    \draw[-{Latex[length=2.0mm]}] (2.95,1.2) -- (4.85,1.2);
    \node[left] at (-0.5,1.45) {$U_1$};
    \node[below] at (3.9,1.12) {$U_2=U_{cb}-U_{db}$};
    \node[anchor=base] at (2.6,-0.68) {\labminorfont 等效分析电路};
  \end{circuitikz}%
}

\newcommand{\TheoryValueTable}{%
  \begin{tabularx}{\linewidth}{@{}Y{0.78}Y{0.90}Y{0.98}Y{0.98}Y{0.98}@{}}
    \toprule
    $R$ & $\omega RC$ & $U_2$理论值 & $\varphi_2$理论值 & $\Delta t$理论值 \\
    \midrule
    $0\Omega$ & 0.000 & $1.000\mathrm{V}$ & $0.00^\circ$ & $0.00\mu\mathrm{s}$ \\
    $30\Omega$ & 0.0886 & $1.000\mathrm{V}$ & $-10.13^\circ$ & $2.81\mu\mathrm{s}$ \\
    $100\Omega$ & 0.2953 & $1.000\mathrm{V}$ & $-32.90^\circ$ & $9.14\mu\mathrm{s}$ \\
    $200\Omega$ & 0.5906 & $1.000\mathrm{V}$ & $-61.13^\circ$ & $16.98\mu\mathrm{s}$ \\
    $300\Omega$ & 0.8859 & $1.000\mathrm{V}$ & $-83.08^\circ$ & $23.08\mu\mathrm{s}$ \\
    $390\Omega$ & 1.1517 & $1.000\mathrm{V}$ & $-98.07^\circ$ & $27.24\mu\mathrm{s}$ \\
    $510\Omega$ & 1.5061 & $1.000\mathrm{V}$ & $-112.83^\circ$ & $31.34\mu\mathrm{s}$ \\
    $1\mathrm{k}\Omega$ & 2.9531 & $1.000\mathrm{V}$ & $-142.58^\circ$ & $39.61\mu\mathrm{s}$ \\
    $2\mathrm{k}\Omega$ & 5.9062 & $1.000\mathrm{V}$ & $-160.78^\circ$ & $44.66\mu\mathrm{s}$ \\
    $5.1\mathrm{k}\Omega$ & 15.0608 & $1.000\mathrm{V}$ & $-172.40^\circ$ & $47.89\mu\mathrm{s}$ \\
    $\infty$ & $\infty$ & $1.000\mathrm{V}$ & $-180.00^\circ$ & $50.00\mu\mathrm{s}$ \\
    \bottomrule
  \end{tabularx}%
}

\newcommand{\MeasurementTable}{%
  \begin{tabularx}{\linewidth}{@{}Y{0.70}Y{0.86}Y{0.86}Y{0.98}Y{0.98}Y{1.62}@{}}
    \toprule
    $R$ & \makecell[c]{$U_1$\\实测/V} & \makecell[c]{$U_2$\\实测/V} &
    \makecell[c]{$\Delta x$\\实测/$\mu$s} & \makecell[c]{$\varphi_2$\\实测} & 备注 \\
    \midrule
    $0\Omega$ & \MeasureBlank & \MeasureBlank & \MeasureBlank & \MeasureBlank & \MeasureBlank \\
    $30\Omega$ & \MeasureBlank & \MeasureBlank & \MeasureBlank & \MeasureBlank & \MeasureBlank \\
    $100\Omega$ & \MeasureBlank & \MeasureBlank & \MeasureBlank & \MeasureBlank & \MeasureBlank \\
    $200\Omega$ & \MeasureBlank & \MeasureBlank & \MeasureBlank & \MeasureBlank & \MeasureBlank \\
    $300\Omega$ & \MeasureBlank & \MeasureBlank & \MeasureBlank & \MeasureBlank & \MeasureBlank \\
    $390\Omega$ & \MeasureBlank & \MeasureBlank & \MeasureBlank & \MeasureBlank & \MeasureBlank \\
    $510\Omega$ & \MeasureBlank & \MeasureBlank & \MeasureBlank & \MeasureBlank & \MeasureBlank \\
    $1\mathrm{k}\Omega$ & \MeasureBlank & \MeasureBlank & \MeasureBlank & \MeasureBlank & \MeasureBlank \\
    $2\mathrm{k}\Omega$ & \MeasureBlank & \MeasureBlank & \MeasureBlank & \MeasureBlank & \MeasureBlank \\
    $5.1\mathrm{k}\Omega$ & \MeasureBlank & \MeasureBlank & \MeasureBlank & \MeasureBlank & \MeasureBlank \\
    $\infty$ & \MeasureBlank & \MeasureBlank & \MeasureBlank & \MeasureBlank & \MeasureBlank \\
    \bottomrule
  \end{tabularx}%
}

\begin{document}
\makelabtitle
\makelabheadertable

\LabMajorSection{移相器调研（20分，不够自行加页）}
\AnalysisParagraph{移相器是在保持信号频率不变的条件下，使输出信号相对输入信号产生指定相位差的电路或系统。常见实现方式包括无源RC移相网络、有源RC移相器、全通滤波器、数字延时移相以及锁相环或直接数字合成移相等。}
\AnalysisParagraph{无源RC移相网络结构简单，适合固定频率或窄频应用，但普通RC串联取样电路的输出幅度会随相位一起变化。若要求输出幅度基本不变，可采用X型RC移相网络或有源全通网络。X型RC网络在理想条件下满足$U_2=U_1$，同时可通过调节$R$实现$0^\circ$至$-180^\circ$的连续滞后相移，符合本实验在$10\mathrm{kHz}$下设计和测试移相器的要求。}
\AnalysisParagraph{有源全通移相器通常由运算放大器和RC网络组成，输入输出隔离较好，带负载能力更强，也便于级联扩展相移范围。数字移相方法可通过采样、延时和重构实现较精确的相位控制，但需要时钟、转换器和数字处理模块。本实验侧重电路原理和参数设计，选取X型RC移相电路作为实现方案。}

\LabMajorSection{实验中移相器的工作原理（10分）}
\LabMinorSection{电路形式}
\begin{center}
  \begin{tabular}{@{}c@{\hspace{0.8cm}}c@{}}
    \resizebox{0.43\linewidth}{!}{\PhaseShiftCircuit} &
    \resizebox{0.43\linewidth}{!}{\EquivalentPhaseCircuit} \\
    {\ExpSixTableFont 图6.1 X型RC移相电路} &
    {\ExpSixTableFont 图6.2 等效分析电路}
  \end{tabular}
\end{center}

\LabMinorSection{相量关系}
设输入信号相量为$\dot U_1=U_1\angle0^\circ$。由等效电路可得
\[
  \dot U_2=\dot U_{cb}-\dot U_{db}
  =\left(\frac{1/(j\omega C)}{R+1/(j\omega C)}
  -\frac{R}{R+1/(j\omega C)}\right)\dot U_1
\]
整理为
\[
  \dot U_2=\frac{1-j\omega RC}{1+j\omega RC}\dot U_1
  =U_1\angle\left[-2\arctan(\omega RC)\right]
\]
因此理想情况下
\[
  U_2=U_1,\qquad \varphi_2=-2\arctan(\omega RC)
\]
当$R=0$时，$\varphi_2=0^\circ$；当$R\to\infty$时，$\varphi_2=-180^\circ$；当$0<R<\infty$时，输出电压相对输入电压的滞后相移在$0^\circ$至$-180^\circ$之间连续变化。

\LabMinorSection{本实验理论值计算}
本实验取$U_1=1.000\mathrm{V}$、$f=10\mathrm{kHz}$、$C=0.047\mu\mathrm{F}$，故
\[
  \omega=2\pi f=6.2832\times10^4\mathrm{rad/s}
\]
各电阻值对应的理论相位按$\varphi_2=-2\arctan(\omega RC)$计算，示波器双通道时间差理论值按
\[
  \Delta t=\frac{|\varphi_2|}{360^\circ f}
\]
计算。
\ExpSixTable{\TheoryValueTable}

\LabMajorSection{实验仪器及器材清单（10分）}
\begin{LabArabicList}
  \item 函数信号发生器：输出$1.000\mathrm{V}$、$10\mathrm{kHz}$正弦信号。
  \item 双踪数字示波器：同时观察输入电压$U_1$和输出电压$U_2$，测量幅值、周期和两通道时间差$\Delta x$。
  \item 数字万用表：校核电阻、电容及连接状态。
  \item 电路实验套件、九孔方板和连接导线。
  \item 电容：$C=0.047\mu\mathrm{F}$，两只。
  \item 电阻：$0\Omega$、$30\Omega$、$100\Omega$、$200\Omega$、$300\Omega$、$390\Omega$、$510\Omega$。
  \item 电阻箱或可变电阻：$1\mathrm{k}\Omega$、$2\mathrm{k}\Omega$、$5.1\mathrm{k}\Omega$及开路等效$\infty$。
\end{LabArabicList}

\LabMajorSection{实验数据与分析（60分）}
\LabMinorSection{实验步骤}
\begin{LabArabicList}
  \item 按图6.1在九孔方板上搭建X型RC移相电路，两个电容均取$0.047\mu\mathrm{F}$。
  \item 将函数信号发生器设置为正弦波，频率$10\mathrm{kHz}$，校准输入电压$U_1=1.000\mathrm{V}$。
  \item 示波器CH1接输入端，CH2接输出端，两个通道共地，记录$U_1$、$U_2$和两波形时间差$\Delta x$。
  \item 依次改变$R$为表中各值，保持$U_1$校准不变，记录各组实测数据。
  \item 由$\varphi_2=-360^\circ f\Delta x$计算实测相位，并与理论值比较。
\end{LabArabicList}

\LabMinorSection{实验数据记录表}
\ExpSixSmallTable{\MeasurementTable}

\LabMinorSection{数据分析}
\AnalysisParagraph{实测完成后，应先检查$U_2$是否接近$U_1$。若$U_2$随$R$明显变化，需检查两个桥臂的电阻、电容是否对称，示波器探头倍率是否一致，以及输出端负载是否过重。}
\AnalysisParagraph{相位分析时，以$10\mathrm{kHz}$对应周期$T=100\mu\mathrm{s}$为基准，由示波器测得的时间差$\Delta x$换算相位：$\varphi_2=-360^\circ\Delta x/T$。负号表示输出电压相对输入电压滞后。}
\AnalysisParagraph{理论上$R$增大时，$\omega RC$增大，$\varphi_2=-2\arctan(\omega RC)$的绝对值单调增大，输出相位由$0^\circ$连续接近$-180^\circ$。若实测曲线保持单调变化且输出幅值基本不变，则说明移相器满足设计要求。}
\LabBlank{3.0\baselineskip}

\LabMajorSection{拓展（10分）}
\AnalysisParagraph{若需实现$0^\circ$至$360^\circ$范围内连续可调，可考虑将两个$0^\circ$至$180^\circ$移相单元级联，或使用有源全通移相器组实现更好的隔离和幅值稳定性。也可在Multisim中建立X型RC网络和有源全通网络模型，对不同$R$值下的幅值、相位和负载影响进行仿真比较。}
\LabBlank{3.0\baselineskip}

\end{document}
